\documentclass[11pt]{article}
\usepackage{amsmath,amssymb,amsthm}
\newtheorem{theorem}{Theorem}\newtheorem{proposition}[theorem]{Proposition}
\title{Retained state, unary invocation recursion, and periodic closure}
\date{16 November 2026}
\begin{document}\maketitle
\paragraph{Audit and interface.} Independently reviewed at the stated unary finite-row tensor-context interface. The target is a precise bridge and obstruction in the inherited native tensor-context, weighted-polynomial doctrine. Retained-state closure stores a state and executes nothing. Invocation recursion generates finite call trees. Instantaneous finite-wire closure selects equal incoming and outgoing wire values. The results below relate explicit semantics built from these operations; they do not identify the operations themselves or discharge the entire feedback-extension contract. In particular, all grades below are full, including unused shapes. No product-coefficient recursive family is identified with a tensor family.

\section{Source pins and parameter restriction}
The retained-state construction used here is \emph{Local Joint Comparison and Generated Interference}, SHA256
\texttt{d6341978238b18be787f9d12158ee7fbc437b144d2d54475a2afe17bf7a833d1},
section ``Feedback closure of the native grade construction.'' It defines machines $(S,f)$, native intertwiners $(h,\alpha)$, retained closure by adjoining state factors, and seeded history prefixes. It explicitly separates retaining state from choosing a seed and executing.
The instantaneous source is \emph{Finite-measure native feedback with retained writer grades}, SHA256
\texttt{74b894e56740f2eec8984a27b9c6467446432f91ace9463638da0cce067ea79b}; its normalized partial-trace domain is \emph{Native Partial-Trace Domains, Staged Elimination, and Shared Coupling}, SHA256
\texttt{9adcdee6ab17ce3071e5c8aa46f0b853cffdfdbd55a6934e5a2fa0bd619a1767}.
The invocation interface is \emph{Initial Row Feedback and Cost Recovery}, SHA256
\texttt{7ffa6f93b05b978f5529f8bf6f2ea79a5d3d34fb4afb1f57dbfa7c44e791da64}, and \emph{Paid Delay and Guarded Source Recovery}, SHA256
\texttt{fdf6961e6f9f9ae6e41a3d425b30e4585b3e7e513b0178ac8fa88cc35535aea9}.

Work with finite normalized raw rows and actual augmented context coefficients in $\mathrm{Vect}_k/k$, tensoring them in chronological order. Content coefficients lie independently in a symmetric monoidal category $\mathcal D$. For terminating semantics use the countable subprobability native target from the arena sources. The bridge is unary: a step either terminates with one result or continues with exactly one state. It does not compile multi-child branching into shared mutable state without additional operational data. A supplied admitted controller may already be included in the state; this is not a theorem that an arbitrary policy descends to a smaller state.

Let $S$ be a state set, $s_0\in S$ an explicitly supplied seed, and $B$ a result set. The step data are
\[
 (W_s,\mu_s,K_{s,w}),\qquad
 d_s:W_s\longrightarrow (B+S)\times O.
\]
The full writer grade is $Q=(O,(D_o))$. Regard it as the unary weighted polynomial $C_Q$, with $(C_Q)_1=O$ and other arities empty. A continuation returns its updated state; a halt returns a $B$ result. Halting is tested after a step. Thus a complete path has length at least one. An initial zero-step halt would be an additional initialization datum and grade summand, not silently included.

\section{Exact native bridge with a necessary grade refinement}
A complete seeded path is
\[
 p=(w_1,\ldots,w_n),\quad s_0\xrightarrow{w_1}s_1\cdots
 \xrightarrow{w_{n-1}}s_{n-1}\xrightarrow{w_n} b,
\]
where the first $n-1$ outputs are continuations and the last is a halt.
It has weight $\prod_{i=1}^n\mu_{s_{i-1}}(w_i)$, context coefficient
$\bigotimes_{i=1}^nK_{s_{i-1},w_i}$, and writer outputs $(o_1,\ldots,o_n)$.
Every raw zero path remains. Let $\mathsf{Term}_{s_0}$ be this independently defined finite-termination semantics with full unary word grade
\[
 \mathcal W(Q)_1=\coprod_{n\geq1}O^n,
 \qquad D_{(o_1,\ldots,o_n)}=\bigotimes_iD_{o_i}.
\]
Its row is countable with mass at most one: from the fixed seed each finite horizon has only finitely many raw paths, by finite row branching, even when $S$ itself is uncountable. The countable union of the finite first-halt layers gives the row; normalized finite-prefix mass conservation bounds its mass by one. Zero paths are included in each such finite family. This is a new, explicit terminal aggregation of state unfoldings, not the assertion that a retained machine already is a terminal process.

Define a stronger full unary grade
\[
 \mathcal W_S(Q)_1=\coprod_{n\geq1}O^n\times S^{n-1},        \tag{1}
\]
with the same content tensor at each labelled shape. The internal state labels are independent formal coordinates, including every word/label combination never attained by any execution. Let $\mathsf{Term}^{\rm lab}_{s_0}$ retain the same raw paths, weights, context coefficients and result, but give $p$ the labelled shape
$((o_1,\ldots,o_n),(s_1,\ldots,s_{n-1}))$.

Define $F:1+S\to B+S$ by using the step at $s_0$ on the external input $*$ and the step at $s$ on a loop input $s$. Its full grade is $C_Q$; each raw witness returns precisely its indicated unary output and writer shape. Let $A(F)$ be its one-paid-event atomic arena, equivalently $DJ(F)$ in the current initial-row doctrine. Its initial states are port-free, so guarded feedback is admitted with no hypothesis on cycles of $S$.

\begin{theorem}[State trajectories as unary invocation recursion]
There is a natural actual native isomorphism
\[
 \operatorname{Eval}_{\rm ind}\operatorname{Fb}_S(A(F))
 \cong \mathsf{Term}^{\rm lab}_{s_0},                    \tag{2}
\]
whose full grade map is the canonical isomorphism
$R_S C_Q\cong\mathcal W_S(Q)$. Each complete invocation node corresponds to exactly one state step, so paid cost is path length. Composing (2) with state-label deletion gives a native comparison
\[
 \operatorname{Eval}_{\rm ind}\operatorname{Fb}_S(A(F))
 \longrightarrow\mathsf{Term}_{s_0}.                    \tag{3}
\]
This comparison is a raw-witness bijection with actual context-coefficient isomorphisms, but its full grade map need not be invertible.
\end{theorem}
\begin{proof}
Every finite formal $R_S C_Q$ tree is a nonempty chain: each node has arity one, and its only slot either ends in an external leaf or carries one $S$-labelled child. Reading node shapes and internal labels from root to leaf gives exactly (1), with inverse the chain constructor. This argument applies to every formal shape, without testing realizability. Its content coefficient comparison is the actual tensor regrouping.

A complete invocation tree is correspondingly a chain of original raw step witnesses. Its incoming labels are forced successively by the preceding continuation outputs. Its first input is $s_0$ and its last output is a halt in $B$, so it gives exactly the displayed seeded path. Conversely every such path constructs that invocation chain. The weight and context coefficient formulas are identical factor by factor, including zero atoms. The terminal result and labelled content body agree. The supplied independent invocation-normal-form theorem for the actual guarded arena therefore gives (2); no unproved contraction of shared scopes is needed. In this unary atomic implementation at most one invocation is active at any time.

The full grade map $\rho_S:\mathcal W_S(Q)\to\mathcal W(Q)$ erases the internal state labels and is identity, up to the stated coherence, on each content tensor. Identity on the retained raw path set, its weights, and context factors satisfies the body square over $\rho_S$, proving (3). No raw histories or coefficients are deleted by this map.
\end{proof}

\paragraph{Native comparison law.}
For a second machine with state $T$, seed $t_0$, row maps $r_s$, full writer map $\gamma:Q\to Q'$, and state map $h:S\to T$, require $h(s_0)=t_0$, exact row pushforward, actual augmented coefficient arrows, and
\[
 d'_{h(s)}(r_s w)=((1_B+h)\times\gamma_O)d_s(w).           \tag{4}
\]
These are the native state-intertwining equations for the tagged stopping interface, not equality of observed transition distributions. Extend $h$ coordinatewise to internal grade labels. Recursive raw path maps preserve first halt and length by (4). At every finite length exact pushforward follows by conditional fibre sums. More explicitly, for a fixed target prefix and each source prefix above it, its current state $s$ satisfies that the target current state is $h(s)$. Summing the source next-step fibre over $r_s$ gives the same prescribed target next-step weight at every such source prefix. Summing the preceding prefix fibre therefore proves the next-prefix equation, without an inverse or section of $h$. The tagged equation (4) reflects halt versus continuation, so restriction to exact first halt preserves this equation. All raw maps are defined also at zero paths; exact fibre equations at zero target atoms have zero total source mass but do not remove their raw atoms or actual coefficient arrows. Countable summation yields the full terminal equation. Tensor the actual coefficient arrows. On all unused grades use the map induced by $\gamma$ on the writer factors and $h$ on the state coordinates. This proves naturality of (2) under these supplied intertwiners, even for noninjective $h$. The squares
\[
 \rho_T\,\mathcal W_h(\gamma)=\mathcal W(\gamma)\,\rho_S
\]
commute on every shape and actual coefficient arrow, so (3) is also natural. A proof only on realized state trajectories would not establish this square.

\section{A strict obstruction to reverse recovery}
If $S$ has at least two elements and $O$ is nonempty, $\rho_S$ identifies different length-two formal shapes. Thus (3) is not an isomorphism of full grades. The following example rules out even an arbitrary reverse native cell between suitable finite cost layers, not just an inverse to this particular map.

\begin{proposition}[Writer-word aggregation cannot recover hidden state grade]
There is a finite normalized unary state machine for which the exact cost-two unlabelled terminal process admits no native arrow to its state-labelled terminal process, although their raw rows and scalar laws are identical.
\end{proposition}
\begin{proof}
Take $S=\{0,1\}$, $s_0=0$, $B=\{b\}$ and $O=\{o\}$. At each state the three raw outcomes are halt to $b$ with weight $1/2$, continue to $0$ with weight $1/4$, and continue to $1$ with weight $1/4$. Every outcome has writer $o$ and unit context/content coefficients. The cost-two terminal layer has exactly two raw atoms, each of mass $1/8$: continue to $0$ then halt, or continue to $1$ then halt. The unlabelled body has the same word $(o,o)$ at both atoms. The labelled body has two distinct shapes $((o,o),0)$ and $((o,o),1)$.

Any native arrow must push its source raw measure to the target raw measure. Both target atoms have positive mass, so both must have preimages. But a single full grade function sends the common source body shape $(o,o)$ to only one target shape. The body equation therefore prevents an image at the other target atom. This contradicts exact pushforward. This is independent of which raw function or actual coefficient arrows are proposed. The full grades remain $\mathcal W(Q)$ and $\mathcal W_S(Q)$, including unused lengths and shapes; those unused source shapes have no raw witnesses that could repair the contradiction.
\end{proof}

The repair is (1), a declared expansion of the retained terminal grade, not a support quotient or an inference that state labels were already present in ordinary writer content. Alternatively an observation may intentionally accept the directed map (3). Retaining state history in the external result can preserve a different observer, but does not by itself make a grade function depend on an individual raw witness. The claimed no-arrow example has a fixed result $b$; it is not a universal no-go for changed output interfaces.

Retained closure alone cannot choose $s_0$. A natural choice of a seed for all nonempty state sets would in particular choose a bit fixed by the bit permutation, which is impossible. A seed-preserving comparison category is genuine additional input. Nor does retained closure specify the tagged halt/continue distinction. Without it there are finite prefixes and possible infinite paths, but not this first-halt terminal row.

\section{Instantaneous closure is exactly a periodic-path restriction}
Now use the finite-measure writer source, let $S$ be finite, and consider a normalized nonstopping body $E:S\to S$ with full writer $Q$, finite raw rows and actual symmetric tensor coefficients. Regard it as a machine with external interface $1\to1$ and exposed state $S$. Write $E^{;n}$ for chronological composition, $n\geq1$. Its full writer grade is $Q^{\boxtimes n}$, not $\mathcal W(Q)$ and not $R_SC_Q$.

For each starting state $s$, let $P_n(s)$ be the complete length-$n$ state trajectories of the body, with product weights, actual step context tensors and writer $n$-tuple. Define a finite-measure raw row
\[
 \mathsf{Per}_n(E)=\{(s,p):s\in S,\ p\in P_n(s),\ s_n(p)=s\}.       \tag{5}
\]
Retain the original path weight, tensor and whole grade $Q^{\boxtimes n}$. This is a separate finite sum/restriction construction; no claim that disjoint union is a categorical coproduct of exact-measure native rows is used. Nor is restriction asserted to be an exact-pushforward native arrow from the unrestricted normalized row.

\begin{theorem}[Periodic trajectory formula, before forgetting native data]
There is an actual native isomorphism
\[
 \operatorname{Tr}_S(E^{;n})\cong \mathsf{Per}_n(E)          \tag{6}
\]
at identity full grade $Q^{\boxtimes n}$, with chronological tensor coherence on context coefficients. The normalized partial trace is defined exactly when
\[
 \sum_{s\in S}\Pr_s[s_n=s]=1.                            \tag{7}
\]
This is fixed-wire naturality and is natural under bijective state readdressing; arbitrary retained-state intertwiners need not induce a closed native comparison.
\end{theorem}
\begin{proof}
A raw witness of $E^{;n}$ at $s$ is its dependent tuple of $n$ consecutive step witnesses. The finite-wire trace retains exactly those tuples whose final state equals $s$, adding $s$ as a raw coordinate. This is literally (5). Its weight, actual context tensor and complete global writer map are the source trace data. Unused writer tuples are retained unchanged. The normalized partial-trace source says precisely that this closed total mass is one, yielding (7), without conditioning or renormalization. A fixed-wire native cell reflects the equality predicate by its body square, so restriction preserves its fibre equation. A bijection on state labels reindexes the seed sum and reflects equality, giving the same conclusion.
\end{proof}

\begin{proposition}[State intertwiners do not generally pass to instantaneous closure]
There is a normalized native retained-state intertwiner whose instantaneous closures have no native comparison in its direction, even with unit coefficients and singleton writer grade.
\end{proposition}
\begin{proof}
Let $E$ deterministically flip the bit state and let $G$ be the deterministic identity on a singleton state. The constant map $h:\{0,1\}\to1$ is an actual state intertwiner: $h\circ\mathrm{flip}=h$, and all raw rows are singleton with weight one and coefficient unit. The full writer map is identity. The instantaneous trace of $E$ has no raw witnesses and mass zero; that of $G$ has one weight-one witness. There is no exact-pushforward native arrow from the former to the latter. In the normalized partial-trace regime the first is undefined and the second defined. Even at horizon two the closed masses are respectively two and one, so the problem is not only a one-step absence of fixed points.
\end{proof}

By contrast the same constant state map preserves every seeded finite unfold, provided the chosen seeds match; its terminal unfolding comparison also exists when a common tagged stopping rule satisfies (4). Summation over all loop inputs and equality filtering introduce precisely the obstruction absent from seeded causal transport. The periodic result is not a seeded execution law: identity on a bit wire has closed mass two, whereas every selected-seed normalized prefix has mass one. For a fresh uniform bit update the closed mass is one, but the retained trace records the coincident initial/final bit, not arbitrary seeded trajectories. Numerical normalization alone cannot turn (6) into recurrence.

\section{Assembly and residuals}
The proved bridge factors through explicit interfaces: retained machine plus seed plus stopping rule gives finite terminal paths; enriching their full grade by all formal internal state labels gives the unary invocation process; forgetting those grade labels gives only a directed native comparison and may have no reverse. Separately, an $n$-step exposed body plus finite-wire equality closure gives periodic paths summed over starting states. These statements supply actual raw and full-grade maps and a counterexample to an otherwise tempting state-intertwiner transport claim.

Nothing here supplies multibranch shared-state recursion, an infinite coefficient tensor, an instantaneous fixed-point selector, arbitrary shared-joint interchange with (6), or an equivalence of unclosed clocks. The usual state prefix tower carries $Q^{\boxtimes n}$ and explicit history outputs; its full comma truncation needs the source's content augmentations. The terminating comparison (2) only tensors completed finite factors and requires no content deletion. Replacing that prefix tower by a fixed full $R_SC_Q$ tower would need the separate partial-terminal-body interface, not a declared shape for an unfinished invocation. These boundaries remain open obligations where required by a larger application.
\end{document}
